\subsection{Beyond electronics: one solver, many physical domains}
\label{sec:domains}

Nothing in the simulator knows about a physical domain. A Verilog-A \texttt{nature} declares a
quantity and its tolerance, a \texttt{discipline} pairs a potential with a flow, and the solver
sees rows of one Newton system whatever they stand for (Section~\ref{sec:arch}). The
examples in Table~\ref{tab:domains} test that claim. Each couples two or more domains through
models written for it in plain Verilog-A, each declares its own natures where the standard
disciplines do not cover it, and each is checked against a result the simulator does not
produce: a closed form, a conservation law, a published figure, or an independent integration
of the same equations. They are demonstrations, not validations against a reference simulator,
because no SPICE can run them. Taken together they use most of the language
(Section~\ref{sec:frontend}): user natures and disciplines, \texttt{analog function}s,
\texttt{generate} loops over vector nets of user disciplines, \texttt{laplace\_nd},
\texttt{transition}, and the event family \texttt{@(cross)}, \texttt{@(timer)},
\texttt{@(above)} and \texttt{@(final\_step)}.
\src{docs/multiphysics.md, its Verilog-A coverage table; docs/photonic-noise.md; docs/traffic.md}

\begin{table*}[!t]
  \centering
  \caption{Examples outside electronics: domains coupled in one system, and the independent check
  each is held to. Every deck is in \texttt{circuits/multiphysics/} (traffic and photonic noise in
  \texttt{circuits/}) and runs as written.}
  \src{docs/multiphysics.md sections 1-15; docs/photonic-noise.md; docs/traffic.md section 6}
  \label{tab:domains}
  \footnotesize
  \setlength{\tabcolsep}{4pt}
  \begin{tabular}{@{}p{0.17\textwidth}p{0.22\textwidth}p{0.29\textwidth}p{0.27\textwidth}@{}}
    \toprule
    Example & Domains & Check & Result \\
    \midrule
    Peltier cooler, PID control & electrical, thermal (absolute K) &
      maximum $\Delta T$ against the closed form; energy balance at an operating point &
      $T_c = 232.4555$\,K vs 232.455\,K; $Q_h - Q_c - P_e = -2.5\times10^{-6}$\,W \\
    Electrostatic MEMS & electrical, mechanical, gas damping &
      pull-in voltage $\sqrt{8kg^3/27\varepsilon A}$; AC resonance and $Q$ at two pressures &
      pull-in at 5.2259\,V as predicted; $Q$ = 1605 and 78.2, matching the models' equations \\
    Filamentary ReRAM & electrical, thermal, filament state &
      SET/RESET voltages under step refinement &
      agree to 0.1\,mV between 1\,\textmu s and 0.2\,\textmu s maximum steps \\
    Hodgkin--Huxley axon & electrical cable, ionic gating, temperature &
      conduction velocity against~\cite{hodgkin1952} &
      18.97\,m/s vs 18.8\,m/s computed in~\cite{hodgkin1952} (0.9\,\%), spatially converged \\
    Power-grid swing & electromechanical (rotors), network &
      equal-area criterion: critical clearing and peak angle &
      peak 130.39$^\circ$ vs 130.40$^\circ$ from the equal areas \\
    Thermo-optic ring bank & optical, thermal, control &
      linewidth and peak temperature against the closed form (16\,ps photon lifetime, 1\,ms thermal
      time, one system) &
      0.1562\,nm vs 0.1564\,nm; 7.1416\,K vs 7.1423\,K; sweep jumps within 0.009\,nm of the folds \\
    Li-ion cell & electrochemical, electrical, thermal &
      capacity at C/20 against the open-circuit curve's closed form; oven abuse test &
      1.9604\,Ah vs 1.9609\,Ah ($-0.03$\,\%); runaway threshold at 134.07\,$^\circ$C \\
    Electromagnetic relay & electrical, magnetic (Jiles--Atherton), mechanical, thermal &
      holding flux density $B = \sqrt{2\mu_0 F/A}$ &
      0.1588\,T vs 0.15853\,T \\
    Gray--Scott reaction--diffusion & chemical (2-D grid, nested \texttt{generate}) &
      well-mixed reactor against its steady states &
      ends at $(0.200000, 0.400000)$, the stable state, to print precision \\
    Piezoelectric harvester & mechanical (Duffing beam), electrical &
      voltage into a matched load against the exact two-degree-of-freedom steady state; Duffing
      jump &
      6.3194\,V vs 6.3193\,V; bistable band 51.762--52.809\,Hz \\
    Circuit-breaker arc & electrical, fluid, thermal (Mayr arc) &
      Mayr static characteristic; recovery voltage after current chopping &
      20\,004\,W vs 20\,000\,W at 10\,A; recovery peak within 1.5\,\% of the closed form \\
    Cavity optomechanics & optical, mechanical, thermal &
      optical damping rate against the eigenvalues of the linearized model; phonon-laser amplitude &
      within 0.1\,\%; 4.7559\,pm vs 4.7577\,pm from an independent integration \\
    Horn loudspeaker & electrical, mechanical, acoustic, thermal &
      input impedance against the exact exponential horn (Webster equation) &
      within 0.011\,\%; radiated power within 0.2\,\% to 1\,kHz at 256 sections \\
    SRAM single-event upset & carrier transport (diffusion), electrical &
      collected charge against the diffusion equation's eigen-series &
      within 0.002\,\% from 2\,ns on; critical charge 10.54\,fC deposited \\
    Microring bistability and chaos & optical, free carriers, thermal &
      self-pulsing period and amplitude against an independent integrator (LSODA, rtol $10^{-10}$) &
      40.500\,ns vs 40.497\,ns; two-ring chaotic peak interval 25.062\,ns vs 25.118\,ns \\
    Photonic phase and detector noise & optical phase, electrical &
      thermal phase-noise PSD~\cite{glenn1989,bartolo2012}; shot and RIN budget of a ring gyroscope &
      within the tabulated PSD's 0.7\,\% interpolation error; budget as computed from the operating point \\
    Motorway ramp metering & traffic (vehicles as a conservative discipline) &
      fundamental diagram, exact; the uncontrolled case of~\cite{bellemans2002} against its own
      discretization &
      fundamental diagram exact; total time spent within 0.6\,\% \\
    \bottomrule
  \end{tabular}
\end{table*}

Building them was also a test of the simulator, and some found what the electronic gates had
not. The circuit-breaker arc exposed a junction-limiting guess that treated the breaker's
terminals as a p-n junction, fixed in 2.16.2. Others were lessons about the models rather than
defects: a microring started exactly on its unstable steady state stays there, as it should, so
its decks switch the laser on with a ramp, as an experiment does. The SRAM upset is the deck that
uses \texttt{.nodeset} to choose one state of a bistable circuit.
\src{docs/multiphysics.md, "What building this found" paragraphs; release.txt 2.16.2}

\subsubsection{A spiking neural network at circuit level}
The largest of these runs is a trained neural network. Xuan \emph{et al.}~\cite{xuan2026snn}
design a spiking neural network as 2T2R memristor crossbars feeding clocked analog
leaky-integrate-and-fire neurons, so that, sampled once per clock cycle, each circuit neuron
matches the software neuron it was trained as. We reproduce it at the level of the paper's
functional diagram: the capacitors (998\,fF and 660\,fF), the RRAM conductances at the stated
levels and the clocked switches are circuit elements, the op-amp, attenuator and comparator are
ideal blocks, and the flip-flops of the refractory timer are built as sample-and-hold circuits.
The neuron's leak is not a parameter but the result of charge sharing between the two
capacitors, $998/(998+660) = 0.602$.

At neuron level, both of the paper's figures are reproduced. In its Fig.~4 (five identical
synapses, 1\,\textmu s refractory period) the circuit spikes at 1.50\,\textmu s and
3.75\,\textmu s, as in the paper. In its Fig.~5 (four synapses from the paper's codebook) it
spikes at steps 6, 10 and 14, as in the paper, with the membrane voltage within 12\,mV of the
paper's equations at every step. Each deck runs in about 3\,s (76 unknowns, about
\num{13000} time points). The paper leaves the attenuator gain open, and its two figures imply
different values (a range of $[2.872, 3.047]$ for Fig.~4 and $[0.656, 0.762]$ for Fig.~5); each
deck uses its own.
\src{circuits/snn/README.md}

At network level, the network was trained on MNIST in software with the circuit's own neuron
model, cycle for cycle, then quantized to the codebook. It reaches 94.04\,\% on \num{10000}
held-out images, against 97.46\,\% reported by the paper, whose input encoding is not specified.
As in the paper, the first hidden layer runs in software. The circuit holds the other three
layers: three crossbars with \num{8832} 2T2R synapses and 138 analog neurons, \num{10973}
devices and \num{2626} unknowns per image. Ten held-out images were run through it. \textbf{The
circuit's class equals the quantized model's on all ten}, and \num{40011} of the \num{40020}
cycle-by-neuron spike slots agree (99.978\,\%). Every disagreement is on an image whose model
membrane passes within about 2\,mV of the threshold. Each image takes 51--64\,min on an
i7-1185G7 under Windows (8 threads; about \num{63000} time points) and 11--14\,min on a
Ryzen~9 9950X under Linux (32 threads), with the same time-point counts, the same spike slots and
the same classes on both. The time follows the number of time points, which the network's
switching activity sets, not the matrix size. The transistor-level version
of the neuron is in progress with IHP's open SG13G2 devices, as the paper's SG13S process with
its memristor module is not public.
\src{circuits/snn/mnist/README.md, results.jsonl, LOG.md step 12; scripts/benchmark/bench\_snn.txt;
release.txt 2.2.0}
